\documentclass[10pt,a4paper]{amsart}
\usepackage{amsmath,amssymb,amsthm,graphicx}
\usepackage{verbatim}
\newenvironment{reference}{\comment}{\endcomment}
%\newenvironment{reference}{}{}
\newenvironment{slogan}{\comment}{\endcomment}
\newenvironment{history}{\comment}{\endcomment}

% For commutative diagrams we use Xy-pic
\usepackage[all]{xy}

% We use 2cell for 2-commutative diagrams.
\xyoption{2cell}
% Theorem environments.
%
\theoremstyle{plain}
\newtheorem{theorem}[subsection]{Theorem}
\newtheorem{proposition}[subsection]{Proposition}
\newtheorem{lemma}[subsection]{Lemma}

\theoremstyle{definition}
\newtheorem{definition}[subsection]{Definition}
\newtheorem{example}[subsection]{Example}
\newtheorem{exercise}[subsection]{Exercise}
\newtheorem{situation}[subsection]{Situation}

\theoremstyle{remark}
\newtheorem{remark}[subsection]{Remark}
\newtheorem{remarks}[subsection]{Remarks}

\numberwithin{equation}{subsection}

% Macros
%
\def\lim{\mathop{\mathrm{lim}}\nolimits}
\def\colim{\mathop{\mathrm{colim}}\nolimits}
\def\Spec{\mathop{\mathrm{Spec}}}
\def\Hom{\mathop{\mathrm{Hom}}\nolimits}
\def\Ext{\mathop{\mathrm{Ext}}\nolimits}
\def\SheafHom{\mathop{\mathcal{H}\!\mathit{om}}\nolimits}
\def\SheafExt{\mathop{\mathcal{E}\!\mathit{xt}}\nolimits}
\def\Sch{\mathit{Sch}}
\def\Mor{\mathop{\mathrm{Mor}}\nolimits}
\def\Ob{\mathop{\mathrm{Ob}}\nolimits}
\def\Sh{\mathop{\mathit{Sh}}\nolimits}
\def\NL{\mathop{N\!L}\nolimits}
\def\CH{\mathop{\mathrm{CH}}\nolimits}
\def\proetale{{pro\text{-}\acute{e}tale}}
\def\etale{{\acute{e}tale}}
\def\QCoh{\mathit{QCoh}}
\def\Ker{\mathop{\mathrm{Ker}}}
\def\Im{\mathop{\mathrm{Im}}}
\def\Coker{\mathop{\mathrm{Coker}}}
\def\Coim{\mathop{\mathrm{Coim}}}

% Boxtimes
%
\DeclareMathSymbol{\boxtimes}{\mathbin}{AMSa}{"02}

%
% Macros for moduli stacks/spaces
%
\def\QCohstack{\mathcal{QC}\!\mathit{oh}}
\def\Cohstack{\mathcal{C}\!\mathit{oh}}
\def\Spacesstack{\mathcal{S}\!\mathit{paces}}
\def\Quotfunctor{\mathrm{Quot}}
\def\Hilbfunctor{\mathrm{Hilb}}
\def\Curvesstack{\mathcal{C}\!\mathit{urves}}
\def\Polarizedstack{\mathcal{P}\!\mathit{olarized}}
\def\Complexesstack{\mathcal{C}\!\mathit{omplexes}}
% \Pic is the operator that assigns to X its picard group, usage \Pic(X)
% \Picardstack_{X/B} denotes the Picard stack of X over B
% \Picardfunctor_{X/B} denotes the Picard functor of X over B
\def\Pic{\mathop{\mathrm{Pic}}\nolimits}
\def\Picardstack{\mathcal{P}\!\mathit{ic}}
\def\Picardfunctor{\mathrm{Pic}}
\def\Deformationcategory{\mathcal{D}\!\mathit{ef}}
\usepackage{cleveref}
\crefname{originalequation}{original equation}{original equations}
\Crefname{originalequation}{Original equation}{Original equations}
\makeatletter
\newlabel{equation-sheafy-trace}{{0B6H}{source}{}{sourcepointer-1}{}}
\newlabel{equation-sheafy-trace@cref}{{[originalequation][0][]0B6H}{[1][0][]source}{}{}{}}
\makeatother
\crefname{originalsection}{original section}{original sections}
\Crefname{originalsection}{Original section}{Original sections}
\makeatletter
\newlabel{cohomology-section-global-RHom}{{0B6A}{source}{}{sourcepointer-2}{}}
\newlabel{cohomology-section-global-RHom@cref}{{[originalsection][0][]0B6A}{[1][0][]source}{}{}{}}
\makeatother
\crefname{originallemma}{original lemma}{original lemmas}
\Crefname{originallemma}{Original lemma}{Original lemmas}
\makeatletter
\newlabel{lemma-iso-on-RSheafHom}{{0A9Q}{source}{}{sourcepointer-3}{}}
\newlabel{lemma-iso-on-RSheafHom@cref}{{[originallemma][0][]0A9Q}{[1][0][]source}{}{}{}}
\makeatother
\title{The trace map induces an isomorphism of global derived Homs}
\author{The Stacks Project Authors}
\date{stacks-project:duality.tex:local-snapshot-20261012; GNU FDL 1.2 or later}
\begin{document}
\maketitle
\noindent\textit{Editorial note.} All delivered mathematical text is the author's own, copied from the pinned Stacks Project chapter file duality.tex; no statement or proof step was rewritten, reordered or removed.\par\medskip
\noindent\textit{Editorial note.} Preamble provenance: the pinned chapter file opens with an input of the shared preamble and carries no preamble text of its own, so the project's single Stacks-wide preamble staged at sources/stacks-duality/preamble.tex is excerpted verbatim under the recipe fixed by LANE\_BRIEF 22.1, as three spans, lines 8-13 (verbatim package and the reference/slogan/history environments), lines 14-18 (Xy-pic setup) and lines 154-217 (the native theorem declarations and mathematical macros). Nothing was added, removed or redefined; the wrapper class and the amsmath/amssymb/amsthm/graphicx base packages are supplied by the assembler.\par\medskip
\noindent\textit{Editorial note.} The counted statement and proof contain no cross reference, no citation and no included file, so the delivered body is self contained.\par\medskip
\noindent\textit{Editorial note.} Canonical location of the counted proposition: Stacks tag 0B6I (sources/stacks-duality/pinned-original-source.tex line 878).\par\medskip
\noindent\textit{Allowed standard prerequisites (editorial).}\begin{itemize}
\item The sheaf level isomorphism underlying the trace map. Same chapter, Lemma labelled lemma-iso-on-RSheafHom, cited at proof lines 920-922. Supplies the isomorphism after applying DQ\_Y.
\item Global derived Hom as derived sections of the internal Hom. Cohomology, Section labelled cohomology-section-global-RHom, cited at proof lines 909-918. Identifies the global Homs with the sections used.
\item The coherator as right adjoint to the inclusion. The adjunction of DQ\_Y used at proof lines 924-931. Shows that DQ\_Y isomorphisms induce isomorphisms on cohomology.
\end{itemize}\medskip
\ifdefined\hypertarget\hypertarget{sourcepointer-1}{}\fi
\noindent\textit{Source reference.} Original equation 0B6H: sources/stacks-duality/pinned-original-source.tex, source lines 758--758. Same-chapter pointer into the author environment equation at source line 758; invoked at proof line None of the counted proof. It is retained only as a navigation pointer; every load-bearing step is inside the delivered span.\par\medskip
\ifdefined\hypertarget\hypertarget{sourcepointer-2}{}\fi
\noindent\textit{Source reference.} Original section 0B6A: sources/stacks-chow/tags.txt, source lines 14063--14063. Cross-chapter standard prerequisite located in the project tag list staged in the local snapshot; invoked at proof line 888 of the counted proof. It is retained only as a navigation pointer; every load-bearing step is inside the delivered span.\par\medskip
\ifdefined\hypertarget\hypertarget{sourcepointer-3}{}\fi
\noindent\textit{Source reference.} Original lemma 0A9Q: sources/stacks-duality/pinned-original-source.tex, source lines 775--775. Same-chapter pointer into the author environment lemma at source line 775; invoked at proof line 899 of the counted proof. It is retained only as a navigation pointer; every load-bearing step is inside the delivered span.\par\medskip
\begin{lemma}
\label{lemma-iso-global-hom}
Let $f : X \to Y$ be a morphism of quasi-separated and quasi-compact
schemes.
For all $L \in D_\QCoh(\mathcal{O}_X)$ and $K \in D_\QCoh(\mathcal{O}_Y)$
(\ref{equation-sheafy-trace}) induces an isomorphism
$R\Hom_X(L, a(K)) \to R\Hom_Y(Rf_*L, K)$ of global derived homs.
\end{lemma}

\begin{proof}
By the construction in
Cohomology, Section \ref{cohomology-section-global-RHom}
we have
$$
R\Hom_X(L, a(K)) =
R\Gamma(X, R\SheafHom_{\mathcal{O}_X}(L, a(K))) =
R\Gamma(Y, Rf_*R\SheafHom_{\mathcal{O}_X}(L, a(K)))
$$
and
$$
R\Hom_Y(Rf_*L, K) = R\Gamma(Y, R\SheafHom_{\mathcal{O}_Y}(Rf_*L, K))
$$
Thus the lemma is a consequence of Lemma \ref{lemma-iso-on-RSheafHom}.
Namely, a map $E \to E'$ in $D(\mathcal{O}_Y)$ which induces
an isomorphism $DQ_Y(E) \to DQ_Y(E')$ induces a quasi-isomorphism
$R\Gamma(Y, E) \to R\Gamma(Y, E')$. Indeed we have
$H^i(Y, E) = \Ext^i_Y(\mathcal{O}_Y, E) = \Hom(\mathcal{O}_Y[-i], E) =
\Hom(\mathcal{O}_Y[-i], DQ_Y(E))$ because $\mathcal{O}_Y[-i]$
is in $D_\QCoh(\mathcal{O}_Y)$ and $DQ_Y$ is the right adjoint
to the inclusion functor $D_\QCoh(\mathcal{O}_Y) \to D(\mathcal{O}_Y)$.
\end{proof}

\end{document}
